\documentclass[12pt]{article}
\usepackage{../../ficha}

\begin{document}

\ficha{Combinatória e Binómio de Newton}[Arranjos, Permutações, Combinações e Triângulo de Pascal]

\instrucoes{Resolve os exercícios, indicando o raciocínio. $^nA_p = \dfrac{n!}{(n - p)!}$, $\;^nC_p = \dfrac{n!}{p!\,(n - p)!}$ e $P_n = n!$.}

\secao{Princípio Fundamental da Contagem}{Resolve.}
\begin{exercicios}
  \item Uma loja vende t-shirts em 4 cores e 5 tamanhos. Quantas t-shirts diferentes há?
  \item Quantas matrículas da forma «dois algarismos, duas letras, dois algarismos» existem (alfabeto de 26 letras)?
  \item Quantos números de três algarismos existem? Quantos têm os algarismos todos diferentes?
  \item Quantos números pares de três algarismos diferentes se podem formar com os algarismos 1, 2, 3, 4 e 5?
  \item Um código tem 4 dígitos (0 a 9). Quantos códigos existem? Quantos não têm dígitos repetidos?
  \item Uma moeda é lançada 6 vezes. Quantas sequências de resultados existem?
\end{exercicios}

\secao{Fatoriais, Arranjos e Permutações}{Calcula ou resolve.}
\begin{exercicios}
  \item $5!$ \; e \; $\dfrac{10!}{8!}$
  \item $^7A_3$ \; e \; $^6A'_2$ \; (arranjos com repetição)
  \item De quantas maneiras se podem sentar 6 pessoas numa fila de 6 cadeiras?
  \item De quantas maneiras se podem atribuir o 1.º, 2.º e 3.º prémios a 10 concorrentes?
  \item Quantos anagramas tem a palavra LIVRO? E a palavra BANANA?
  \item De quantas maneiras se podem sentar 5 pessoas numa fila, se duas delas querem ficar juntas?
  \item De quantas maneiras se podem sentar 6 pessoas à volta de uma mesa redonda?
  \item Resolve: $\dfrac{(n + 1)!}{(n - 1)!} = 30$.
\end{exercicios}

\secao{Combinações}{Resolve.}
\begin{exercicios}
  \item Calcula $^8C_3$ e $^{10}C_{7}$.
  \item De quantas maneiras se pode escolher uma comissão de 3 pessoas de um grupo de 12?
  \item Num totoloto escolhem-se 6 números de 1 a 49. Quantas chaves diferentes existem?
  \item Quantas diagonais tem um octógono?
  \item Uma turma tem 12 rapazes e 10 raparigas. Quantas comissões de 4 alunos com 2 rapazes e 2 raparigas se podem formar?
  \item Na mesma turma, quantas comissões de 4 alunos têm pelo menos uma rapariga?
  \item Num baralho de 52 cartas, quantas mãos de 5 cartas têm exatamente 2 ases?
  \item Quantos caminhos mínimos há numa grelha de $5 \times 3$ quadrículas, do canto inferior esquerdo ao canto superior direito?
\end{exercicios}

\secao{Triângulo de Pascal}{Resolve.}
\begin{exercicios}
  \item Escreve as linhas 0 a 7 do triângulo de Pascal.
  \item Qual é a soma dos elementos da linha 10?
  \item Numa linha do triângulo de Pascal, o segundo elemento é 15. Quantos elementos tem a linha? Qual é o maior?
  \item Mostra que $^nC_p = {}^nC_{n-p}$ e explica o que isso significa no triângulo.
  \item Verifica que $^6C_2 + {}^6C_3 = {}^7C_3$ e enuncia a propriedade geral.
  \item Numa linha do triângulo de Pascal, a soma dos elementos é 256. Que linha é?
\end{exercicios}

\secao{Binómio de Newton}{Resolve.}
\begin{exercicios}
  \item Desenvolve $(a + b)^4$.
  \item Desenvolve $(x - 2)^5$.
  \item Desenvolve $(2x + 1)^3$.
  \item Determina o 4.º termo do desenvolvimento de $(x + 3)^6$ (ordenado por potências decrescentes de $x$).
  \item Determina o coeficiente de $x^3$ no desenvolvimento de $(2x - 1)^7$.
  \item Determina o termo independente de $x$ em $\left(x + \dfrac{1}{x}\right)^{8}$.
  \item Determina o termo independente de $x$ em $\left(x^2 - \dfrac{2}{x}\right)^{6}$.
  \item Calcula $1{,}01^5$ com quatro casas decimais, usando $(1 + 0{,}01)^5$.
\end{exercicios}

\secao{Desafios}{Resolve os problemas.}
\begin{exercicios}
  \item Quantos números de 4 algarismos têm os algarismos em ordem estritamente crescente?
  \item De quantas maneiras se podem distribuir 10 bolas iguais por 3 caixas diferentes? (Sugestão: «estrelas e barras».)
  \item Mostra que $\sum_{p=0}^{n} {}^nC_p = 2^n$ usando o binómio de Newton.
  \item Numa festa, cada pessoa cumprimentou todas as outras uma vez. Houve 66 cumprimentos. Quantas pessoas estavam na festa?
\end{exercicios}

\end{document}
